\documentclass[12pt,a4paper]{article}
\usepackage[utf8]{vietnam}
\usepackage{amsmath,amssymb}
\usepackage{geometry}
\geometry{
  top=12mm,
  bottom=12mm,
  left=14mm,
  right=14mm,
  footskip=6mm,
  headheight=14pt
}

\usepackage{tikz}
\usetikzlibrary{arrows.meta,calc,shapes.geometric}
\usepackage[most]{tcolorbox}
\usepackage{fontawesome}
\usepackage{fancyhdr}
\usepackage{enumitem}

\pagestyle{fancy}
\fancyhf{}
\fancyhead[L]{\footnotesize\color{teal!80!black}\faBook\ \textbf{SỔ TAY KIẾN THỨC TOÁN 3}}
\fancyhead[R]{\footnotesize\color{teal!80!black}\textbf{Lớp Toán Thầy Lĩnh -- loptoanthaylinh.vn}}
\fancyfoot[C]{\footnotesize\color{gray}-- Sổ tay kiến thức: Bài 03. Tìm thành phần trong phép cộng, phép trừ --}
\renewcommand{\headrulewidth}{0.4pt}
\renewcommand{\footrulewidth}{0pt}

\setlength{\parindent}{0pt}
\setlength{\parskip}{0pt}

\begin{document}

% --- BANNER BÀI HỌC ---
\begin{tcolorbox}[
  enhanced,
  colback=teal!8!white,
  colframe=teal!70!black,
  arc=3mm, boxrule=1.2pt,
  left=3mm, right=3mm, top=2mm, bottom=2mm,
  halign=center
]
  {\LARGE\bfseries\color{teal!90!black}\faBookmark\ BÀI 03: TÌM THÀNH PHẦN TRONG PHÉP CỘNG, PHÉP TRỪ}\\[1mm]
  {\itshape\footnotesize\color{black!75} Quy tắc tìm số hạng chưa biết, tìm số bị trừ và tìm số trừ trên 1 trang A4}
\end{tcolorbox}

\vspace{-2mm}

\noindent
\begin{minipage}[t]{0.49\textwidth}\vspace{0pt}
\begin{tcolorbox}[
  enhanced,
  colback=blue!4!white,
  colframe=blue!75!black,
  arc=2mm, boxrule=0.8pt,
  left=2mm, right=2mm, top=1.8mm, bottom=1.8mm,
  title={\footnotesize\bfseries\color{white}\faKey\ 1. Tìm Số hạng chưa biết}
]
{\footnotesize
\begin{itemize}\setlength{\itemsep}{1pt}\setlength{\leftskip}{-2.5mm}
  \item Cấu trúc: $\mathbf{\text{Số hạng}} + \mathbf{\text{Số hạng}} = \mathbf{\text{Tổng}}$.
  \item \textbf{Quy tắc vàng:} Muốn tìm một số hạng chưa biết, ta lấy \textbf{Tổng trừ đi số hạng đã biết}:
  $$x + a = b \implies x = b - a$$
  $$a + x = b \implies x = b - a$$
  \item \textit{Bản chất:} Tổng luôn là giá trị lớn nhất trong phép cộng.
\end{itemize}
}
\end{tcolorbox}
\end{minipage}
\hfill
\begin{minipage}[t]{0.49\textwidth}\vspace{0pt}
\begin{tcolorbox}[
  enhanced,
  colback=teal!4!white,
  colframe=teal!75!black,
  arc=2mm, boxrule=0.8pt,
  left=2mm, right=2mm, top=1.8mm, bottom=1.8mm,
  title={\footnotesize\bfseries\color{white}\faSliders\ 2. Tìm Số bị trừ \& Số trừ}
]
{\footnotesize
\begin{itemize}\setlength{\itemsep}{1pt}\setlength{\leftskip}{-2.5mm}
  \item Cấu trúc: $\mathbf{\text{Số bị trừ}} - \mathbf{\text{Số trừ}} = \mathbf{\text{Hiệu}}$.
  \item \textbf{Tìm Số bị trừ (số đứng đầu - lớn nhất):}
  $$\text{Số bị trừ} = \text{Hiệu} + \text{Số trừ}$$
  $$x - a = b \implies x = b + a$$
  \item \textbf{Tìm Số trừ (số ở giữa):}
  $$\text{Số trừ} = \text{Số bị trừ} - \text{Hiệu}$$
  $$a - x = b \implies x = a - b$$
\end{itemize}
}
\end{tcolorbox}
\end{minipage}

\vspace{-2.5mm}
% --- KHỐI 3: SƠ ĐỒ TƯ DUY TIKZ ---
\noindent
\begin{tcolorbox}[
  enhanced,
  colback=orange!4!white,
  colframe=orange!80!black,
  arc=2mm, boxrule=0.8pt,
  left=2mm, right=2mm, top=1.5mm, bottom=1.5mm,
  title={\footnotesize\bfseries\color{white}\faSitemap\ SƠ ĐỒ TƯ DUY: TAM GIÁC THẦN KÌ TÌM $x$ CỘNG - TRỪ}
]
\centering
\begin{tikzpicture}[
  box/.style={draw, rounded corners=2pt, align=center, font=\scriptsize\bfseries, inner sep=3pt, minimum height=0.85cm},
  arr/.style={-{Stealth[length=2.5mm]}, thick}
]
  \node[box, fill=blue!12, draw=blue!80!black, text width=3.3cm] (k1) at (0, 0) {TÌM SỐ HẠNG ($x$)\\$x + a = b$\\$\mathbf{x = b - a}$ (Lấy Tổng trừ)};
  \node[box, fill=purple!12, draw=purple!80!black, text width=3.3cm] (k2) at (4.7, 0) {TÌM SỐ BỊ TRỪ ($x$)\\$x - a = b$\\$\mathbf{x = b + a}$ (Lấy Hiệu cộng)};
  \node[box, fill=teal!12, draw=teal!80!black, text width=3.3cm] (k3) at (9.4, 0) {TÌM SỐ TRỪ ($x$)\\$a - x = b$\\$\mathbf{x = a - b}$ (Lấy SBT trừ)};

  \draw[arr, blue!80!black] (k1) -- (k2);
  \draw[arr, purple!80!black] (k2) -- (k3);
\end{tikzpicture}
\end{tcolorbox}

\vspace{-2.5mm}
% --- KHỐI 4: VÍ DỤ & MẸO TRÁNH BẪY ---
\noindent
\begin{minipage}[t]{0.49\textwidth}\vspace{0pt}
\begin{tcolorbox}[
  enhanced,
  colback=green!4!white,
  colframe=green!60!black,
  arc=2mm, boxrule=0.8pt,
  left=2mm, right=2mm, top=1.8mm, bottom=1.8mm,
  title={\footnotesize\bfseries\color{white}\faCheckCircle\ Ví dụ thực hành chuẩn mực}
]
{\footnotesize
\textbf{1. Tìm số hạng:}\\
$x + 245 = 680 \implies x = 680 - 245 = \mathbf{435}$.\\
\textbf{2. Tìm số bị trừ (đứng đầu):}\\
$x - 136 = 254 \implies x = 254 + 136 = \mathbf{390}$.\\
\textbf{3. Tìm số trừ (ở giữa):}\\
$512 - x = 180 \implies x = 512 - 180 = \mathbf{332}$.
}
\end{tcolorbox}
\end{minipage}
\hfill
\begin{minipage}[t]{0.49\textwidth}\vspace{0pt}
\begin{tcolorbox}[
  enhanced,
  colback=red!4!white,
  colframe=red!70!black,
  arc=2mm, boxrule=0.8pt,
  left=2mm, right=2mm, top=1.8mm, bottom=1.8mm,
  title={\footnotesize\bfseries\color{white}\faExclamationTriangle\ Mẹo tránh bẫy trắc nghiệm}
]
{\footnotesize
\begin{itemize}\setlength{\itemsep}{1.5pt}\setlength{\leftskip}{-2.5mm}
  \item \textbf{Bẫy nhầm Số bị trừ và Số trừ:}\\
  $x - a = b$ thì $x = b + a$ (phép CỘNG).\\
  $a - x = b$ thì $x = a - b$ (phép TRỪ).\\
  Nhiều bạn cứ thấy dấu trừ là auto làm phép cộng!
  \item \textbf{Mẹo thử lại:}\\
  Sau khi tìm xong $x$, hãy thay $x$ trở lại đề bài ban đầu để tính nhẩm xem có đúng kết quả không.
\end{itemize}
}
\end{tcolorbox}
\end{minipage}


\end{document}
